\section{Discussion}
\label{sec:discussion}

\subsection{Principal findings}
On more than 99{,}000 real inpatient encounters, a boosting-anchored
knowledge-gated sequence model (TKGN-B) was the best or joint-best model for
both 30-day readmission and next-stay treatment escalation under
patient-disjoint validation, was clearly and significantly better than
three established deep sequence architectures, and was the best model on the
prospective temporal split for treatment escalation. Its advantage over
carefully tuned gradient boosting was, however, small (about 0.002 AUROC on
average) and not statistically significant on a single test set, and on the
temporal readmission split a random forest was more robust than every other
model. The most important ingredient of the sequence model was access to
the patient's earlier encounters (and, for readmission, their outcomes); the
knowledge-graph components did not improve discrimination, neither with the
full training set nor when training data were scarce.

\subsection{Interpretation}
\paragraph{Why anchoring helps}
The tabular learners already receive hand-crafted history aggregates, so
the remaining signal that a sequence model can add is limited. Training the
network on the residual of an out-of-fold boosting model concentrates its
capacity on this remaining signal: the order and content of earlier stays,
changes between stays, and patients whose history is long enough to
saturate simple counts (Table~\ref{tab:subgroups}). The same residual
network trained on its own (TKGN) could not match boosting, which agrees
with the broader observation that neural networks struggle to beat trees on
heterogeneous tabular data \citep{grinsztajn2022}. Combining both
through an additive logit, rather than choosing between them, is a simple
way to keep the strengths of each; the out-of-fold construction is
essential, because in-sample boosting logits on the training set are
optimistically sharp and would leave the network nothing to learn.

\paragraph{The role of clinical knowledge}
Ontology-based code representations were originally motivated by rare codes
and small datasets \citep{choi2017gram}. In this cohort, most diagnosis
codes are stored at the three-character level, and the frequent codes that
carry most of the signal have hundreds to thousands of training examples,
which leaves little room for hierarchical smoothing. The empirical
co-morbidity propagation even reduced performance slightly, possibly because
PPMI emphasises strongly associated but rare pairs (for example
pregnancy-related categories) that are not predictive of the outcomes
studied here. The dedicated data-scarcity experiment
(Section~\ref{sec:kgeff}) showed that this is not only a large-sample
effect: even with 5--25\% of the training patients the graph did not help.
We think the hierarchy is more likely to matter in settings with
full-length ICD codes, many more distinct codes per patient, smaller
institutions or transfer between sites, and we report the negative result
to discourage over-claiming of knowledge-graph benefits on this
benchmark.

\paragraph{History-reliability gates}
The gates did not change discrimination in the ablation, but they made the
model's use of history inspectable. For escalation they opened steadily with
the length of the history, which matches the intuition that a longer
medication trajectory is more informative for the next treatment decision.
For readmission they damped the history summary while the explicit outcome
tokens of earlier stays remained important, suggesting that ``what happened
after the last discharge'' is more useful than a compressed summary of
earlier stays.

\paragraph{Temporal robustness}
The temporal split revealed shifts in recording practice (more diagnoses per
stay) and in the population (longer histories, more emergency visits) that
are typical of real hospital data. Such shifts are invisible to random
splits and strongly penalised models that extrapolate linearly. They also
explain why model rankings differed between designs. Any deployment of
these models would need local recalibration and monitoring of feature
distributions over time \citep{vancalster2019}.

\paragraph{Fairness}
Discrimination was similar for women and men and for the two largest racial
groups, and calibration was acceptable in all adequately sized subgroups.
The larger gap across age bands, with the weakest performance in the oldest
patients, mirrors the higher and more homogeneous baseline risk of this
group and deserves attention because older adults are a key target
population for transitional-care programmes. Because the race variable is
coarse and partly missing, these results do not rule out inequities for
smaller groups \citep{obermeyer2019}.

\subsection{Comparison with previous work}
The discrimination reached for 30-day readmission (AUROC around 0.69) lies
at the upper end of the range typically reported for readmission models
\citep{kansagara2011}. Direct comparison with earlier analyses of the same
database is difficult because many use encounter-level random splits, which
allow repeat stays of a patient to appear in training and test data, or
restrict the cohort to one encounter per patient. We deliberately used
patient-disjoint and temporal validation, multiple repeats and validation-only
tuning, which give more conservative but more trustworthy estimates.
Treatment escalation at the next hospital stay has, to our knowledge, not
been used as a prediction target on this database; it links the readmission
literature with work on clinical inertia in type 2 diabetes
\citep{khunti2013}.

\subsection{Limitations}
This study has several limitations. First, only one public data source was
available; although the temporal split mimics prospective use,
external validation in an independent health system is needed. Second, the
release has no calendar dates, so stays were ordered by encounter identifier
and the time between stays is unknown; time-aware models could exploit
this information where it exists. Third, only three diagnoses per stay are
recorded and most codes are truncated to three characters, which limits the
potential of ontology-based representations. Fourth, readmissions are only
observed within the participating systems, so some outcomes may be missed.
Fifth, treatment escalation is defined from prescription states at the next
inpatient stay and is therefore a proxy for glycaemic deterioration rather
than a direct measure of it; it is only defined for patients who return to
hospital. Sixth, the learning curves are based on a single repeat per
training size, and the differences between the best models are small
relative to their variability; claims of superiority over gradient boosting
should therefore be made with caution. Seventh, the neural architectures
were developed on a separate partition of the same cohort, so a small
optimistic bias from design choices cannot be excluded; the tabular
baselines were tuned only on the validation partitions of the evaluation
runs. Finally, the models are research
prototypes and have not been evaluated for clinical utility.

\subsection{Conclusions}
Using real, publicly available diabetes inpatient records and a
leakage-controlled protocol, we showed that anchoring a knowledge-gated
sequence model on out-of-fold gradient boosting yields a well-calibrated
model that matched or slightly exceeded tuned gradient boosting under
patient-disjoint validation and was clearly better than standard deep
sequence models, for both 30-day readmission and next-stay treatment
escalation; under temporal validation it was the best model for escalation
but not for readmission. Most of the gain from
sequence modelling comes from the patient's earlier stays, whereas
knowledge-graph code sharing did not improve discrimination at any training
size examined. All code, derived results and an
interactive dashboard are released to support replication and extension to
other cohorts.
